\documentclass[11pt]{amsart}
\pagestyle{plain}
\raggedbottom
\usepackage[dvipsnames]{xcolor}
\usepackage[utf8]{inputenc}
\usepackage{amssymb,amsmath,amsthm,amsfonts,fullpage,float,latexsym,bbm,microtype,cite,fancyvrb}
\usepackage{enumerate}
\usepackage{hyperref}
\hypersetup{colorlinks=true, linkcolor=blue, citecolor=magenta, filecolor=magenta, urlcolor=magenta}
\usepackage{bm,mathrsfs}

\newtheorem{theorem}{Theorem}
\newtheorem{proposition}[theorem]{Proposition}
\newtheorem{claim}[theorem]{Claim}
\newtheorem{conjecture}[theorem]{Conjecture}
\newtheorem{lemma}[theorem]{Lemma}
\newtheorem{corollary}[theorem]{Corollary}
\theoremstyle{definition}
\newtheorem{remark}[theorem]{Remark}
\newtheorem{definition}[theorem]{Definition}
\newtheorem{observation}[theorem]{Observation}
\newtheorem{example}[theorem]{Example}

\renewcommand{\emptyset}{\varnothing}
\newcommand{\C}{\mathbb{C}}
\newcommand{\Z}{\mathbb{Z}}
\DeclareMathOperator{\Hilb}{Hilb}
\DeclareMathOperator{\OSP}{OSP}
\DeclareMathOperator{\inv}{inv}
\DeclareMathOperator{\Inv}{Inv}
\DeclareMathOperator{\sgn}{sgn}
\DeclareMathOperator{\blocks}{blocks}
\DeclareMathOperator{\GL}{GL}
\DeclareMathOperator{\Sym}{Sym}
\DeclareRobustCommand{\qbinom}{\genfrac[]{0pt}{}}

\begin{document}

\title{A combinatorial interpretation of the hook coefficients of the Hilbert
series of the multigraded coinvariant ring}
\maketitle

\section{The Problem}

Let $m$ and $n$ be positive integers and let $R_n^{(m)}$ be the coinvariant
algebra with $m$ sets of $n$ variables. Its multigraded Hilbert series admits a
Schur expansion
$$
\Hilb(R_n^{(m)};q_1,\ldots ,q_m)= \sum_{\lambda}
c_{\lambda}(n)\,s_\lambda(q_1,\ldots ,q_m),
$$
with nonnegative coefficients $c_\lambda(n)$ that do not depend on $m$. We give a
combinatorial interpretation of $c_\lambda(n)$ when $\lambda$ is a \emph{hook
shape} partition, i.e.\ $\lambda=(a,1^b)$ with $a\ge1,\ b\ge0$, or
$\lambda=\varnothing$.

\medskip
\noindent\textbf{Answer (Theorem~\ref{thm:improved-coefficient-formula}).}
For $a \geq 1$ and $b \geq 0$,
$$
c_{(a,1^b)}(n)
= \#\{w \in \OSP(n,n-b)\ :\ \inv(w) = a,\ w \text{ is of type I} \},
$$
where $\OSP(n,n-b)$ is the set of ordered set partitions of $\{1,\dots,n\}$ into
$n-b$ blocks, $\inv$ is the inversion statistic of Sagan--Swanson, and ``type I''
is the mergeability condition recalled in Section~\ref{sec:phi}. Moreover
$c_\varnothing(n)=1$ and $c_\lambda(n)=0$ for every non-hook $\lambda$.

\section{Background}

\subsection{The two coinvariant rings}
The \textbf{coinvariant ring with $m$ sets of variables} is
\begin{equation}\label{eq:m-coinv}
R_n^{(m)} = \C[x_i^{(j)}:1\le i\le n,\ 1\le j\le m]
/\big\langle \C[x_i^{(j)}]_+^{\mathfrak{S}_n} \big\rangle,
\end{equation}
where the subscript $+$ denotes the polynomials with vanishing constant term and
the superscript $\mathfrak{S}_n$ denotes invariance under the diagonal action of
$\mathfrak{S}_n$ permuting the $n$ variables within each of the $m$ sets. It is
$\Z_{\ge0}^m$-multigraded, the $j$-th degree counting the total degree in the
$j$-th set $x_1^{(j)},\dots,x_n^{(j)}$. The associated multigraded Hilbert series
is
\begin{equation}
\Hilb(R_n^{(m)}; q_1,\ldots,q_m) :=
\sum_{r_1,\ldots,r_m \geq 0}\dim\big((R_n^{(m)})_{r_1,\ldots,r_m}\big)
q_1^{r_1}\cdots q_m^{r_m}.
\end{equation}

The \textbf{superspace coinvariant ring} is
\begin{equation}\label{eq:super-coinv}
R_n^{(1,1)} = \C[x_1,\ldots,x_n, \theta_1, \ldots,\theta_n]
/\big\langle \C[x_1,\ldots,x_n, \theta_1, \ldots,\theta_n]_+^{\mathfrak{S}_n}
\big\rangle,
\end{equation}
where the $x_i$ are commuting and the $\theta_i$ are anticommuting
($\theta_i\theta_j=-\theta_j\theta_i$, so $\theta_i^2=0$), and $\mathfrak{S}_n$
acts diagonally. It is bigraded: each $x_i$ contributes to the $q$-degree and
each $\theta_i$ to the $u$-degree. Its bigraded Hilbert series is
\begin{equation}
\Hilb(R_n^{(1,1)}; q;u) :=
\sum_{r,s\ge0}\dim\big((R_n^{(1,1)})_{r,s}\big)q^{r}u^{s}.
\end{equation}

\subsection{Symmetric-function tools}
Let $[k]_q := 1+q+\cdots+q^{k-1}$, $[k]_q!:=[k]_q[k-1]_q\cdots[1]_q$ with
$[0]_q!:=1$, and define the $q$-Stirling numbers of the second kind $S[n,k]$ by
$S[0,k]=\delta_{0,k}$, $S[n,0]=\delta_{n,0}$, and
\begin{equation}\label{eq:stirling}
S[n,k] = S[n-1,k-1] + [k]_q\,S[n-1,k]\qquad (0\le k\le n).
\end{equation}
We write $s_\lambda$ for a Schur function and $s_{\mu/\nu}$ for a skew Schur
function. The (single super-alphabet) \textbf{super Schur function} in one
commuting variable $q$ and one anticommuting variable $u$ is
\begin{equation}\label{eq:superschur-def}
s_\lambda(q/u) = \sum_{\nu \subseteq \lambda}
s_\nu(q)\,s_{\lambda'/\nu'}(u),
\end{equation}
where $\lambda'$ is the conjugate (transpose) partition and the sum is over
partitions $\nu$ contained in $\lambda$
(see \cite[Sections A.2.2 and 3.2]{ChengWangBook}).

\subsection{Ordered set partitions and the inversion statistic}
An \textbf{ordered set partition} of $[n]:=\{1,\dots,n\}$ into $k$ blocks is a
sequence $w=(S_1/S_2/\cdots/S_k)$ of pairwise disjoint nonempty sets with
$S_1\cup\cdots\cup S_k=[n]$. We write $\OSP(n,k)$ for the set of these, and
$\blocks(w):=k$. Following Sagan--Swanson \cite[Section 5.1]{SaganSwanson2024},
the inversion set and statistic are
\begin{equation}
\Inv(w) = \{(s,S_j)\ :\ s \in S_i \text{ for some } i < j
\text{ and } s \geq \min S_j\},\qquad \inv(w):=\#\Inv(w).
\end{equation}

\section{Reduction to the superspace coinvariant ring}\label{sec:reduction}

The hook coefficients $c_{(a,1^b)}(n)$ are defined through the \emph{commuting}
multigraded ring $R_n^{(m)}$ of \eqref{eq:m-coinv}, but our combinatorial formula
will be read off the \emph{super} ring $R_n^{(1,1)}$ of \eqref{eq:super-coinv}.
This section makes the bridge precise. It rests on two external theorems
--- a representation-stability theorem of Bergeron \cite{Bergeron2020} and the
diagonal supersymmetry theorem of Lentfer \cite{Lentfer-Supersymmetry} --- which
we invoke rather than reprove. Accordingly we (a) state each theorem in the form
in which it is proved, (b) verify that its hypotheses are met by our rings, and
(c) check the combinatorial identification $P(1,1,n)=\{\text{hooks of length}\le n\}$
that makes the two coefficient sequences interchangeable. Throughout, ``the same
$c_\lambda(n)$'' means the literally identical integer sequence, indexed by
partitions $\lambda$ and depending only on $n$; the whole point of the reduction
is that diagonal supersymmetry identifies the two a priori different sequences.

\subsection{The general multigraded super coinvariant ring}
To state both theorems uniformly we recall the family that contains both
$R_n^{(m)}$ and $R_n^{(1,1)}$. For $\ell,m\ge0$ let
\begin{equation}\label{eq:super-mixed}
R_n^{(\ell,m)}
= \C[\,x_i^{(1)},\dots,x_i^{(\ell)},\ \theta_i^{(1)},\dots,\theta_i^{(m)}: 1\le i\le n\,]
\big/\big\langle (\cdots)_+^{\mathfrak{S}_n}\big\rangle,
\end{equation}
where the $\ell$ sets $x^{(1)},\dots,x^{(\ell)}$ are commuting and the $m$ sets
$\theta^{(1)},\dots,\theta^{(m)}$ are anticommuting, $\mathfrak{S}_n$ acts
diagonally (simultaneously permuting the subscript $i$ in every set), and the
ideal is generated by the $\mathfrak{S}_n$-invariants of positive total degree.
Then $R_n^{(m)}=R_n^{(m,0)}$ (no anticommuting variables) and
$R_n^{(1,1)}$ is the case $\ell=m=1$. The ring $R_n^{(\ell,m)}$ is multigraded by
the $\ell+m$ degrees, the commuting degrees recorded by $q_1,\dots,q_\ell$ and the
anticommuting degrees by $u_1,\dots,u_m$.

\subsection{Bergeron's representation-stability theorem}
The classical mechanism (polarization of coinvariants) is that the $\ell$ sets of
commuting variables span an $\ell$-dimensional space on which $\GL_\ell(\C)$ acts,
and this action commutes with the diagonal $\mathfrak{S}_n$-action. Hence
$R_n^{(\ell,0)}=R_n^{(\ell)}$ is a polynomial $\GL_\ell\times\mathfrak{S}_n$-module
in which the $\GL_\ell$-weight is exactly the commuting multidegree. We use this in
the following form.

\begin{theorem}[Bergeron \cite{Bergeron2020}]\label{prop:bergeron}
Fix $n\ge1$. There is a sequence of nonnegative integers $\big(c_\lambda(n)\big)$,
indexed by partitions $\lambda$ and depending only on $n$ (not on $m$), with
$c_\lambda(n)=0$ whenever $\ell(\lambda)>n$, such that for every $m\ge1$
\begin{equation}\label{eq:bergeron-expansion}
\Hilb\big(R_n^{(m)}; q_1,\ldots,q_m\big)
= \sum_{\lambda:\ \ell(\lambda)\le n} c_\lambda(n)\, s_\lambda(q_1,\ldots,q_m).
\end{equation}
Concretely, $c_\lambda(n)$ is the multiplicity of the irreducible
$\GL$-representation indexed by $\lambda$ in the polarized coinvariant module, equal
to the multiplicity of the $\mathfrak{S}_n$-irreducible component in degree
$\lambda$; this multiplicity is intrinsic to $n$ and $\lambda$ and does not
reference $m$.
\end{theorem}

\begin{proof}[Hypothesis check and invocation]
This is the multigraded Schur-positivity and stability statement established in
\cite{Bergeron2020} (it is the unmixed, purely bosonic specialization of the
results there; the existence of the universal Schur expansion is also the form in
which the problem statement quotes it). We verify the three hypotheses on which it
relies and explain why each conclusion we use follows.

\smallskip\noindent\emph{(1) The relevant module is a polynomial
$\GL_m\times\mathfrak{S}_n$-module.} The variables $x_i^{(j)}$ ($1\le j\le m$) for
fixed $i$ span a copy of the defining representation $V=\C^m$ of $\GL_m$, and
$R_n^{(m)}=R_n^{(m,0)}$ is built from $\Sym(V\otimes\C^n)$ by quotienting by the
diagonally $\mathfrak{S}_n$-invariant ideal. The $\GL_m$-action (linear
substitutions among the $m$ sets) commutes with the diagonal $\mathfrak{S}_n$-action
because the latter permutes the subscript $i$ uniformly across all $m$ sets, while
the former acts on the superscript $j$. The defining ideal is generated by
$\mathfrak{S}_n$-invariants of positive degree, which form a $\GL_m$-stable set
(a $\GL_m$-substitution of $\mathfrak{S}_n$-invariants is again
$\mathfrak{S}_n$-invariant), so the quotient is a $\GL_m$-module. Thus $R_n^{(m)}$
is a polynomial $\GL_m$-module, and its $\GL_m$-character --- which by definition
of the weight grading is exactly $\Hilb(R_n^{(m)};q_1,\dots,q_m)$ --- is a
symmetric polynomial with a Schur expansion. Schur positivity
($c_\lambda(n)\ge0$) is then the statement that these are honest representation
multiplicities; finiteness in each multidegree (so the multiplicities are finite
nonnegative integers) holds because $R_n^{(m)}$ is a finite-dimensional quotient
of a polynomial ring by an ideal generated in positive degree.

\smallskip\noindent\emph{(2) $m$-independence.} The coefficient $c_\lambda(n)$ is
the multiplicity of the $\GL_m$-irreducible of highest weight $\lambda$, which by
$(\GL_m,\GL_\infty)$ branching / the polarization principle equals the
$\mathfrak{S}_n$-isotypic multiplicity attached to the partition $\lambda$ in the
coinvariant module; this is an invariant of the pair $(n,\lambda)$ alone. In
particular, raising $m$ to $m'>m$ adjoins more variables but does not change the
multiplicity of any $\lambda$ with $\ell(\lambda)\le m$: passing from $m$ to $m+1$
sets is the standard stabilization of $\GL$-characters, under which Schur functions
$s_\lambda(q_1,\dots,q_m)$ extend to $s_\lambda(q_1,\dots,q_{m+1})$ with the same
coefficient. This is exactly the assertion in the problem statement that
``$c_\lambda(n)$ \dots\ do not depend on $m$,'' and is what \cite{Bergeron2020}
proves.

\smallskip\noindent\emph{(3) Length bound and saturation.} A Schur polynomial
$s_\lambda(q_1,\dots,q_m)$ is identically $0$ when $\ell(\lambda)>m$ and is a
nonzero polynomial when $\ell(\lambda)\le m$. Hence in \eqref{eq:bergeron-expansion}
only $\lambda$ with $\ell(\lambda)\le m$ contribute. Moreover $c_\lambda(n)=0$
whenever $\ell(\lambda)>n$: the coinvariant module is generated by at most $n$
``polarization directions'' (the $\mathfrak{S}_n$-isotypic types are indexed by
partitions of size $\le n$, so no $\GL$-highest weight with more than $n$ rows can
occur), an integral part of the polarization theorem in \cite{Bergeron2020}.
Combining, every coefficient $c_\lambda(n)$ that is ever nonzero has
$\ell(\lambda)\le n$, and for such $\lambda$ the function $s_\lambda(q_1,\dots,q_m)$
is already nonzero once $m\ge n$. Therefore \emph{all} nonzero coefficients are
already visible in the single instance $m=n$, and conversely no new coefficients
appear for $m>n$. This is the precise meaning of ``all of the possible
$c_\lambda(n)$ that could appear do appear once $m\ge n$,'' so the determination of
the sequence $\big(c_\lambda(n)\big)$ may be carried out at any single $m\ge n$, or
indeed in any setting whose Hilbert series is governed by the same sequence.
\end{proof}

\begin{remark}\label{rem:saturation}
Theorem~\ref{prop:bergeron} reduces ``compute $c_\lambda(n)$ for all $m$'' to
``compute the single $m$-free sequence $c_\lambda(n)$.'' Because the sequence is
$m$-free, we are free to read it off from \emph{any} generating function in which
it appears with linearly independent coefficient functions. Diagonal supersymmetry
(Theorem~\ref{prop:ds}) provides exactly such a generating function for the hook
indices, in the bigraded ring $R_n^{(1,1)}$.
\end{remark}

\subsection{The index set $P(1,1,n)$}
Let $P(\ell,m,n)$ denote the set of partitions $\lambda$ with $\lambda_{\ell+1}\le m$
and $\ell(\lambda)\le n$. For $\ell=m=1$ this is the set we need.

\begin{lemma}\label{lem:P11hooks}
$P(1,1,n)=\{\varnothing\}\cup\{(a,1^b): a\ge1,\ b\ge0,\ b+1\le n\}$, which is
exactly the set of hook partitions of length at most $n$.
\end{lemma}

\begin{proof}
By definition $\lambda\in P(1,1,n)$ iff $\lambda_2\le1$ and $\ell(\lambda)\le n$.
The condition $\lambda_2\le1$ says all parts after the first are $\le1$: either
$\lambda=\varnothing$, or $\lambda_1=a\ge1$ and the remaining parts are a string of
$1$'s, i.e.\ $\lambda=(a,1^b)$ with $b\ge0$. These are precisely the hook shapes
(arm $a$, leg $b$). For $\lambda=(a,1^b)$ the constraint $\ell(\lambda)\le n$ reads
$b+1\le n$; for $\lambda=\varnothing$ it is vacuous. Thus $P(1,1,n)$ is exactly the
set of hooks of length $\le n$ (together with $\varnothing$, the empty hook).
\end{proof}

\subsection{Diagonal supersymmetry and instantiation at $(1,1)$}
Diagonal supersymmetry relates the purely commuting expansion of
$R_n^{(m,0)}=R_n^{(m)}$ to the mixed expansion of $R_n^{(\ell,m)}$: the same
universal coefficients $c_\lambda(n)$ govern both, with the ordinary Schur
functions replaced by \emph{super} Schur functions in the commuting alphabet
$q_1,\dots,q_\ell$ and the anticommuting alphabet $u_1,\dots,u_m$. We record the
theorem and then specialize $(\ell,m)=(1,1)$.

\begin{theorem}[Diagonal supersymmetry; conjectured in \cite{Bergeron2020},
proved in \cite{Lentfer-Supersymmetry}]\label{prop:ds-general}
Fix $n\ge1$ and let $\big(c_\lambda(n)\big)$ be the universal sequence of
Theorem~\ref{prop:bergeron}. For all $\ell,m\ge1$,
\begin{equation}\label{eq:ds-general}
\Hilb\big(R_n^{(\ell,m)};q_1,\dots,q_\ell;u_1,\dots,u_m\big)
= \sum_{\lambda\in P(\ell,m,n)} c_\lambda(n)\,
s_\lambda(q_1,\dots,q_\ell/u_1,\dots,u_m),
\end{equation}
where $s_\lambda(q_1,\dots,q_\ell/u_1,\dots,u_m)$ is the super Schur function of
\eqref{eq:superschur-def} (extended to several variables), and the sum is over
$\lambda$ with $\lambda_{\ell+1}\le m$ and $\ell(\lambda)\le n$.
\end{theorem}

The content of \cite{Lentfer-Supersymmetry} is precisely that the sequence
$c_\lambda(n)$ occurring in \eqref{eq:ds-general} is \emph{the same} sequence as in
\eqref{eq:bergeron-expansion} --- it does not introduce a new, a priori different
set of constants. We now specialize.

\begin{proposition}[Instantiation at $(\ell,m)=(1,1)$]\label{prop:ds}
Fix $n\ge1$. With the universal coefficients $c_\lambda(n)$ of
Theorem~\ref{prop:bergeron},
\begin{equation}\label{eq:coeffs-Hilbert}
\Hilb\big(R_n^{(1,1)}; q;u\big) = \sum_{\lambda \in P(1,1,n)}
c_{\lambda}(n)\, s_\lambda(q/u).
\end{equation}
\end{proposition}

\begin{proof}[Hypothesis check and invocation]
We verify that $R_n^{(1,1)}$ is the instance $(\ell,m)=(1,1)$ of $R_n^{(\ell,m)}$
in \eqref{eq:super-mixed}, and apply Theorem~\ref{prop:ds-general} there.

\smallskip\noindent\emph{(1) $R_n^{(1,1)}$ is the $(1,1)$ instance.} Setting
$\ell=m=1$ in \eqref{eq:super-mixed} gives one set of commuting variables
$x_1,\dots,x_n$ and one set of anticommuting variables $\theta_1,\dots,\theta_n$,
with the diagonal $\mathfrak{S}_n$-action and the ideal generated by
positive-degree $\mathfrak{S}_n$-invariants. This is verbatim the definition
\eqref{eq:super-coinv} of $R_n^{(1,1)}$. The commuting degree is recorded by
$q=q_1$ and the anticommuting degree by $u=u_1$, matching the bigrading of
$R_n^{(1,1)}$. Thus the hypotheses of Theorem~\ref{prop:ds-general} hold with
$\ell=m=1$.

\smallskip\noindent\emph{(2) Specializing \eqref{eq:ds-general}.} With $\ell=m=1$,
the index set is $P(1,1,n)$ and the super Schur functions are
$s_\lambda(q/u)$. Theorem~\ref{prop:ds-general} therefore gives exactly
\eqref{eq:coeffs-Hilbert}, with the \emph{same} $c_\lambda(n)$ as in
\eqref{eq:bergeron-expansion}; this identity of coefficient sequences is the
diagonal supersymmetry theorem and is what guards against any definitional drift
between the commuting and super settings.

\smallskip\noindent\emph{(3) Only hooks survive, and they are all captured.} By
Lemma~\ref{lem:schur-1-1} below, $s_\lambda(q/u)=0$ unless $\lambda$ is a hook, and
the surviving indices are exactly the hooks, which by Lemma~\ref{lem:P11hooks} are
exactly the elements of $P(1,1,n)$. Hence the right-hand side of
\eqref{eq:coeffs-Hilbert} involves precisely the hook coefficients
$c_{(a,1^b)}(n)$ (and $c_\varnothing(n)$), each multiplied by a super Schur function
$s_{(a,1^b)}(q/u)=q^au^b+q^{a-1}u^{b+1}$ whose monomials are determined by
$(a,b)$. Because these $2$-term polynomials are linearly independent enough to be
solved for the individual $c_{(a,1^b)}(n)$ (carried out explicitly in
Lemma~\ref{lem:c_{(a,1^b)}}), the bigraded series of $R_n^{(1,1)}$ determines
\emph{every} hook coefficient of the universal sequence. By
Theorem~\ref{prop:bergeron} this sequence is $m$-independent, so these are the hook
coefficients of $R_n^{(m)}$ for all $m$ at once.
\end{proof}

Thus determining the hook coefficients $c_{(a,1^b)}(n)$ for all $m$ is equivalent
to extracting them from \eqref{eq:coeffs-Hilbert}, which we do in
Sections~\ref{sec:superschur}--\ref{sec:phi}.

\section{The hook super Schur function}\label{sec:superschur}

\begin{lemma}\label{lem:schur-1-1}
Let $\lambda$ be a partition.
\begin{enumerate}[\rm(i)]
\item $s_{\varnothing}(q/u)=1$.
\item If $\lambda=(a,1^b)$ with $a\ge1$ and $b\ge0$, then
$s_\lambda(q/u)= q^a u^b + q^{a-1}u^{b+1}$.
\item If $\lambda$ is not a hook (equivalently $\lambda_2\ge2$), then
$s_\lambda(q/u)=0$.
\end{enumerate}
\end{lemma}

\begin{proof}
We evaluate \eqref{eq:superschur-def}, recalling two single-variable facts.

\emph{Fact A.} In one commuting variable $q$, $s_\nu(q)=0$ unless $\nu$ is a
single row $\nu=(r)$ (including $r=0$, i.e.\ $\nu=\varnothing$), in which case
$s_{(r)}(q)=q^{r}$. Indeed $s_\nu(q)$ is the generating function of column-strict
fillings of $\nu$ with entries in $\{1\}$; such a filling exists iff every column
has length $\le1$, i.e.\ $\nu$ has a single row, and then there is exactly one
filling, of content $r$.

\emph{Fact B.} In one (commuting) variable $u$, the skew Schur function
$s_{\mu/\nu}(u)$ equals $u^{|\mu/\nu|}$ if $\mu/\nu$ is a horizontal strip
($\nu_i\le\mu_i$ and $\mu_{i+1}\le\nu_i$ for all $i$), and is $0$ otherwise; this
is the count of semistandard fillings of $\mu/\nu$ with entries in $\{1\}$, which
is $1$ exactly when $\mu/\nu$ is a horizontal strip.

In \eqref{eq:superschur-def} the second factor is $s_{\lambda'/\nu'}(u)$, a skew
Schur in the single variable $u$. By Fact B it is nonzero iff $\lambda'/\nu'$ is a
horizontal strip, equivalently $\lambda/\nu$ is a \emph{vertical} strip (conjugate
of a horizontal strip). Combining with Fact A: a term survives in
\eqref{eq:superschur-def} iff
\begin{equation}\label{eq:survival}
\nu=(r)\ \text{is a single row, \quad and\quad } \lambda/\nu \text{ is a vertical strip.}
\end{equation}

\emph{(i)} For $\lambda=\varnothing$ the only $\nu$ is $\varnothing=(0)$, with
$s_\varnothing(q)=1$ and $s_{\varnothing/\varnothing}(u)=1$, giving
$s_\varnothing(q/u)=1$.

\emph{(ii)} Let $\lambda=(a,1^b)$, $a\ge1$. A single-row $\nu=(r)$ satisfies
$\nu\subseteq\lambda$ iff $0\le r\le a$. The skew shape $\lambda/(r)$ has
$a-r$ cells in row $1$ and one cell in each of rows $2,\dots,b+1$. This is a
vertical strip (at most one cell per row) iff $a-r\le1$, i.e.\ $r\in\{a-1,a\}$
(and $r=a-1$ requires $a\ge1$, which holds).
\begin{itemize}
\item $r=a$: $\nu=(a)$, $|\lambda/\nu|=b$, contributing
$q^{a}\cdot u^{b}$.
\item $r=a-1$: $\nu=(a-1)$, $|\lambda/\nu|=b+1$, contributing
$q^{a-1}\cdot u^{b+1}$.
\end{itemize}
For each surviving $\nu$ the corresponding $\lambda/\nu$ is indeed a vertical
strip, so its single-variable skew Schur is $u^{|\lambda/\nu|}$ by Fact B. Hence
$s_{(a,1^b)}(q/u)=q^a u^b+q^{a-1}u^{b+1}$.

\emph{(iii)} Suppose $\lambda$ is not a hook, so $\lambda_2\ge2$. For any
single-row $\nu=(r)$ we have $\nu_i=0$ for $i\ge2$, while $\lambda$ has
$\lambda_2\ge2$ cells in its second row; hence row $2$ of $\lambda/\nu$ contains
$\lambda_2\ge2$ cells, so $\lambda/\nu$ is not a vertical strip. By
\eqref{eq:survival} no term survives and $s_\lambda(q/u)=0$.
\end{proof}

\section{The alternating formula for the hook coefficients}\label{sec:alternating}

We will use the following enumeration of Sagan--Swanson.

\begin{proposition}[\!\!{\cite[Theorem 5.1]{SaganSwanson2024}}]\label{prop:OSP-inv}
For $n,k\ge0$,
\begin{equation}\label{eq:OSP-equation}
[k]_q!\,S[n,k] = \sum_{w \in \OSP(n,k)}q^{\inv(w)},
\qquad\text{i.e.}\qquad
\langle q^i\rangle\,[k]_q!\,S[n,k]=\#\{w\in\OSP(n,k):\inv(w)=i\}.
\end{equation}
\end{proposition}

We also use the Rhoades--Wilson Hilbert series.

\begin{theorem}[\!\!\cite{RhoadesWilson2023}]\label{thm:hilbert}
For $n\ge1$,
\begin{equation}\label{eq:hilbert-rhoades-wilson}
\Hilb(R_n^{(1,1)};q;u) = \sum_{k=0}^n [k]_q!\,S[n,k]\,u^{\,n-k}.
\end{equation}
\end{theorem}

\begin{lemma}\label{lem:c_{(a,1^b)}}
For $a\ge1$ and $b\ge0$,
\begin{equation}\label{eq:first-formula-coeffs}
c_{(a,1^b)}(n) =
\sum_{i=0}^b (-1)^i\,\#\{w \in \OSP(n,n-b+i)\ :\ \inv(w) = a+i \}.
\end{equation}
\end{lemma}

\begin{proof}
Equating \eqref{eq:coeffs-Hilbert} and \eqref{eq:hilbert-rhoades-wilson} and
applying Lemma~\ref{lem:schur-1-1} to the left-hand side (only hooks and
$\varnothing$ contribute, by Lemma~\ref{lem:P11hooks}), we get
\begin{equation}\label{eq:master}
1 + \sum_{\substack{a \geq 1\\ b \geq 0}} c_{(a,1^b)}(n)\,
\big(q^a u^b+q^{a-1}u^{b+1}\big)
= \sum_{d=0}^n [d]_q!\,S[n,d]\,u^{n-d}.
\end{equation}
By \eqref{eq:OSP-equation} the right-hand side equals
\begin{equation}\label{eq:rhs-osp}
\sum_{d=0}^n \sum_{i\ge0} \#\{w\in\OSP(n,d):\inv(w)=i\}\,q^i\,u^{n-d}.
\end{equation}

We now extract the coefficient of $q^a u^b$ on both sides
($a\ge1$, $b\ge0$). On the left, a monomial $q^au^b$ arises from exactly two
hook terms: the $q^{a}u^{b}$ term of $(a,1^b)$ (present since $a\ge1$, $b\ge0$),
and, when $b\ge1$, the $q^{(a+1)-1}u^{(b-1)+1}=q^au^b$ term of
$(a+1,1^{b-1})$. No other hook contributes: the $(\alpha,1^\beta)$ term
$q^{\alpha}u^{\beta}+q^{\alpha-1}u^{\beta+1}$ equals $q^au^b$ only if
$(\alpha,\beta)=(a,b)$ (first monomial) or $(\alpha,\beta)=(a+1,b-1)$ (second
monomial); the constant $1$ contributes only to $q^0u^0$ and is irrelevant since
$a\ge1$. On the right, the coefficient of $q^au^b$ in \eqref{eq:rhs-osp} comes
from the stratum $d=n-b$ (so that $u^{n-d}=u^{b}$) and $i=a$.

Hence, for $b\ge1$,
\begin{equation}\label{eq:recurrence}
c_{(a,1^b)}(n)+c_{(a+1,1^{b-1})}(n)
= \#\{w\in\OSP(n,n-b):\inv(w)=a\},
\end{equation}
and, for $b=0$,
\begin{equation}\label{eq:base}
c_{(a)}(n) = \#\{w\in\OSP(n,n):\inv(w)=a\}.
\end{equation}

We solve the recurrence by telescoping in the index $b$. Define, for $i\ge0$,
\begin{equation}
A_i := \#\{w\in\OSP(n,n-b+i):\inv(w)=a+i\}.
\end{equation}
We claim $c_{(a,1^b)}(n)=\sum_{i=0}^b(-1)^iA_i$. Apply
\eqref{eq:recurrence} with the substitutions $(a,b)\mapsto(a+i,b-i)$, valid for
$0\le i\le b-1$ since then $b-i\ge1$:
\begin{equation}\label{eq:tele}
c_{(a+i,1^{b-i})}(n)
= \#\{w\in\OSP(n,n-(b-i)):\inv(w)=a+i\}
- c_{(a+i+1,1^{b-i-1})}(n)
= A_i - c_{(a+i+1,1^{b-i-1})}(n).
\end{equation}
Starting from $i=0$ and iterating \eqref{eq:tele} for $i=0,1,\dots,b-1$,
\begin{align*}
c_{(a,1^b)}(n)
&= A_0 - c_{(a+1,1^{b-1})}(n)
= A_0 - A_1 + c_{(a+2,1^{b-2})}(n)
= \cdots\\
&= \sum_{i=0}^{b-1}(-1)^i A_i + (-1)^b c_{(a+b,1^0)}(n).
\end{align*}
By the base case \eqref{eq:base},
$c_{(a+b)}(n)=\#\{w\in\OSP(n,n):\inv(w)=a+b\}=A_b$ (taking $i=b$ above:
$\OSP(n,n-b+b)=\OSP(n,n)$ and $\inv=a+b$). Therefore
$c_{(a,1^b)}(n)=\sum_{i=0}^b(-1)^iA_i$, which is
\eqref{eq:first-formula-coeffs}.
\end{proof}

\section{The Sagan--Swanson involution and a manifestly positive formula}
\label{sec:phi}

\subsection{The involution $\phi$}
We recall the involution $\phi$ of Sagan--Swanson
\cite[Section 5.1]{SaganSwanson2024}. Given $w=(S_1/\cdots/S_k)\in\OSP(n,k)$ and
$i\in\{1,\dots,k\}$, write $M=\max S_i$. We say $w$ is:
\begin{itemize}
\item \textbf{splittable} at $M=\max S_i$ if $|S_i|>1$;
\item \textbf{mergeable} at $M=\max S_i$ (for $i\le k-1$) if $|S_i|=1$ and
$M>\max S_{i+1}$.
\end{itemize}
Find the largest $M$ at which $w$ is splittable or mergeable, if one exists. Then
$w$ is of \textbf{type I} if it is mergeable at this $M$; \textbf{type II} if it
is splittable at this $M$; \textbf{type III} if no such $M$ exists. Define
\begin{equation}\label{eq:phi}
\phi(w)=\begin{cases}
(S_1/\cdots/S_{i-1}/S_i\cup S_{i+1}/S_{i+2}/\cdots/S_k)
& \text{type I (merge at $S_i$)},\\[2pt]
(S_1/\cdots/S_{i-1}/\{M\}/S_i\setminus\{M\}/S_{i+1}/\cdots/S_k)
& \text{type II (split $S_i$ at $M$)},\\[2pt]
w & \text{type III}.
\end{cases}
\end{equation}

We first record a convenient reformulation of $\inv$ and a structural fact about
the active value $M$, both proved directly from the definition of $\Inv$.

\begin{lemma}[Crossing form of $\inv$]\label{lem:inv-cross}
For $w=(S_1/\cdots/S_k)$ and blocks $S_a,S_b$ put
$\operatorname{cr}(S_a,S_b):=\#\{s\in S_a: s\ge\min S_b\}$. Then
\begin{equation}\label{eq:inv-cross}
\inv(w)=\sum_{1\le a<b\le k}\operatorname{cr}(S_a,S_b).
\end{equation}
\end{lemma}

\begin{proof}
By definition $\Inv(w)=\{(s,S_b): s\in S_a\text{ for some }a<b,\ s\ge\min S_b\}$.
Since the blocks are pairwise disjoint, each element $s$ lies in a unique block
$S_a$; hence the pair $(s,S_b)$ determines the index $a$ and the constraint $a<b$.
Thus $\Inv(w)$ is the disjoint union over ordered index pairs $a<b$ of the sets
$\{(s,S_b): s\in S_a,\ s\ge\min S_b\}$, the $(a,b)$-th of which has cardinality
$\operatorname{cr}(S_a,S_b)$. Summing gives \eqref{eq:inv-cross}.
\end{proof}

\begin{lemma}[Position of large values]\label{lem:large-values}
Let $w=(S_1/\cdots/S_k)$ be not of type III, and let $M=\max S_i$ be the largest
value at which $w$ is splittable or mergeable (so $M$ is the value used by
$\phi$). Then every value $v>M$ lies in a block $S_p$ with $p\ge i$, and
moreover such a block is a singleton $S_p=\{v\}$. Consequently no value $\ge M$
occurs in any of the blocks $S_1,\dots,S_{i-1}$.
\end{lemma}

\begin{proof}
Let $v>M$ and let $S_p$ be the block containing $v$.

\emph{$S_p$ is a singleton.} If $|S_p|\ge2$ then $w$ is splittable at
$\max S_p\ge v>M$, contradicting the maximality of $M$. Hence $S_p=\{v\}$.

\emph{If $p<k$ then $v<\max S_{p+1}$.} The values are distinct, so either
$v<\max S_{p+1}$ or $v>\max S_{p+1}$. In the latter case $S_p=\{v\}$ is a
singleton with $v>\max S_{p+1}$, i.e.\ $w$ is mergeable at $v>M$, again
contradicting maximality of $M$. Hence $v<\max S_{p+1}$, and in particular
$\max S_{p+1}>v>M$.

We now show $p\ge i$. Suppose for contradiction that $p<i$. Apply the two facts
above repeatedly: from $\max S_{p+1}>M$ (just shown) we get that $S_{p+1}$
contains a value $>M$, so by the same two facts $S_{p+1}$ is a singleton and, if
$p+1<k$, $\max S_{p+2}>\max S_{p+1}>M$. Iterating, for every index $q$ with
$p\le q\le k$ we obtain $\max S_q>M$, and each such $S_q$ (for $q<k$) is a
singleton with strictly increasing maxima. In particular this applies to $q=i$
(since $p<i\le k$), giving $\max S_i>M$. But $\max S_i=M$ by hypothesis, a
contradiction. Therefore $p\ge i$.

Finally, since the block $S_i$ contains $M$ and all values are distinct, no block
$S_1,\dots,S_{i-1}$ contains $M$; and by the paragraph above no such block
contains any value $>M$ either. Hence no value $\ge M$ occurs in
$S_1,\dots,S_{i-1}$.
\end{proof}

\begin{lemma}\label{lem:phi-props}
$\phi$ is an involution on $\bigcup_{k}\OSP(n,k)$. If $w$ is of type I then
$\phi(w)$ is of type II with $\blocks(\phi(w))=\blocks(w)-1$ and
$\inv(\phi(w))=\inv(w)-1$; if $w$ is of type II then $\phi(w)$ is of type I with
$\blocks(\phi(w))=\blocks(w)+1$ and $\inv(\phi(w))=\inv(w)+1$. Type III elements
are the fixed points. Consequently $n-\blocks(w)-\inv(w)$ and the parity of
$\blocks(w)$ behave as stated, and the quantity $\inv(w)+\blocks(w)$ increases by
$2$ under a split and decreases by $2$ under a merge.
\end{lemma}

\begin{proof}
\emph{Involution and types.} For a type-I $w$, the merge at $S_i=\{M\}$ produces
$\phi(w)$ in which the block at position $i$ is $S_i\cup S_{i+1}\ni M$ with
$|S_i\cup S_{i+1}|\ge2$; we claim $M$ is again the largest splittable/mergeable
value of $\phi(w)$ and that $\phi(w)$ is splittable there, so $\phi$ is type II at
the same $M$ and $\phi(\phi(w))=w$. Indeed, the merge does not change any value's
position relative to the largest values: by Lemma~\ref{lem:large-values} all
values $>M$ already lay in singleton blocks of index $\ge i+1$ in $w$
(equivalently $\ge i+1$ in $\phi(w)$, whose blocks past position $i$ are exactly
$S_{i+2},\dots$), so $\phi(w)$ is neither splittable nor mergeable at any value
$>M$; and it is splittable at $M=\max(S_i\cup S_{i+1})$ since that block has size
$\ge2$. Splitting $S_i\cup S_{i+1}$ at $M$ removes $M$ as a singleton in front,
recovering $w$. The type-II case is the mirror image: for splittable $w$ with
$M=\max S_i$, $|S_i|\ge2$, Lemma~\ref{lem:large-values} shows all values $>M$ are
singleton blocks of index $\ge i+1$, so after splitting off $\{M\}$ the new
singleton $\{M\}$ at position $i$ is mergeable (its successor $S_i\setminus\{M\}$
has maximum $<M$, as $M=\max S_i$) and $M$ remains the largest active value; thus
$\phi(w)$ is type I and merging undoes the split. Type III elements have no $M$,
so $\phi$ fixes them; these are exactly the fixed points. Hence $\phi$ is an
involution interchanging types I and II at the same value $M$.

\emph{Effect on $\blocks$.} A merge replaces two blocks by their union
($\blocks\mapsto\blocks-1$); a split replaces one block by two
($\blocks\mapsto\blocks+1$). This is immediate from \eqref{eq:phi}.

\emph{Effect on $\inv$ (the load-bearing computation).} It suffices to treat the
type-I merge; the split is its inverse, so $\Delta\inv$ for a split is $+1$ once
we show $\Delta\inv=-1$ for the corresponding merge. Let $w$ be type I, merging
$S_i=\{M\}$ into $B:=S_{i+1}$ (so $M>\max B\ge\min B$, whence $\min B<M$). Write
the blocks of $w$ as
\[
\underbrace{S_1,\dots,S_{i-1}}_{\text{left }P},\ \{M\},\ B,\
\underbrace{S_{i+2},\dots,S_k}_{\text{right }Q},
\]
and of $\phi(w)$ as $P,\ D,\ Q$ with $D:=\{M\}\cup B$. Note
\begin{equation}\label{eq:minD}
\min D=\min(\{M\}\cup B)=\min B\qquad(\text{since }M>\min B).
\end{equation}
We compare $\inv(w)$ and $\inv(\phi(w))$ using the crossing form of Lemma~\ref{lem:inv-cross}
(equation~\eqref{eq:inv-cross}), grouping the ordered block pairs by the blocks involved.

\emph{Pairs internal to $P$, internal to $Q$, and $P$–$Q$ pairs} are identical in
$w$ and $\phi(w)$ (the blocks of $P$ and of $Q$, and their left-to-right order,
are unchanged), so they contribute equally.

\emph{Pairs with left part in $P$ and right part the merged block.} In $w$ these
are, for each $S_a\in P$, the two terms $\operatorname{cr}(S_a,\{M\})$ and
$\operatorname{cr}(S_a,B)$; in $\phi(w)$ the single term
$\operatorname{cr}(S_a,D)$. By \eqref{eq:minD},
$\operatorname{cr}(S_a,D)=\#\{s\in S_a:s\ge\min B\}=\operatorname{cr}(S_a,B)$.
Hence the change contributed by this group, summed over $S_a\in P$, is
\[
\sum_{S_a\in P}\Big(\operatorname{cr}(S_a,D)-\operatorname{cr}(S_a,\{M\})-\operatorname{cr}(S_a,B)\Big)
=-\sum_{S_a\in P}\operatorname{cr}(S_a,\{M\})
=-\#\{s\in S_1\cup\cdots\cup S_{i-1}: s\ge M\}.
\]
By Lemma~\ref{lem:large-values}, no value $\ge M$ occurs in $S_1,\dots,S_{i-1}$,
so this sum is $0$: the $P$–merged-block pairs contribute the same total in $w$
and $\phi(w)$.

\emph{Pairs with left part the merged block and right part in $Q$.} For each
$S_c\in Q$, in $w$ these are $\operatorname{cr}(\{M\},S_c)$ and
$\operatorname{cr}(B,S_c)$; in $\phi(w)$ the single term
$\operatorname{cr}(D,S_c)$. Since $D=\{M\}\sqcup B$ is a disjoint union,
$\operatorname{cr}(D,S_c)=\#\{s\in D:s\ge\min S_c\}
=\operatorname{cr}(\{M\},S_c)+\operatorname{cr}(B,S_c)$. Hence this group
contributes equally in $w$ and $\phi(w)$.

\emph{The pair $(\{M\},B)$.} This ordered pair exists in $w$ (positions $i<i+1$)
but has no analogue in $\phi(w)$, where $M$ and $B$ are in the same block. Its
contribution is $\operatorname{cr}(\{M\},B)=\#\{s\in\{M\}:s\ge\min B\}=1$, because
$M>\max B\ge\min B$.

Summing the four groups, every group is unchanged except the last, which is
present in $w$ and absent in $\phi(w)$. Therefore
\[
\inv(\phi(w))-\inv(w)=-\operatorname{cr}(\{M\},B)=-1.
\]
Thus a merge lowers $\inv$ by exactly $1$, and dually a split raises $\inv$ by
exactly $1$.

\emph{Conclusion.} Combining, type I has $\blocks\mapsto\blocks-1$,
$\inv\mapsto\inv-1$, and type II has $\blocks\mapsto\blocks+1$,
$\inv\mapsto\inv+1$. In particular $n-\blocks(w)-\inv(w)$ is preserved, the parity
of $\blocks(w)$ flips on non-fixed points, and $\inv(w)+\blocks(w)$ changes by
$\pm2$, as stated. (These recover the corresponding assertions of
\cite[Section 5.1]{SaganSwanson2024}, here derived from the definition of
$\Inv$.)
\end{proof}

\subsection{The signed set and its cancellation}
Fix integers $\ell\ge1$ and $0\le b\le\lfloor(\ell-1)/2\rfloor$. Define
\begin{equation}\label{eq:Sdef}
S_{n,\ell,b}=\bigcup_{c=0}^b \{w\in\OSP(n,n-c):\inv(w)=\ell-b-c\},
\end{equation}
a disjoint union over $c$ (the strata have distinct block counts $n-c$). On
$S_{n,\ell,b}$ we use the sign
\begin{equation}\label{eq:epsdef}
\varepsilon(w):=(-1)^{\,\blocks(w)-(n-b)}=(-1)^{\,b-c}\qquad
\text{for }w\in\OSP(n,n-c).
\end{equation}

\begin{lemma}[Standing window]\label{lem:window}
If $w\in S_{n,\ell,b}$ then $n-b\le\blocks(w)\le n$ and
$\ell-2b\le\inv(w)\le\ell-b$.
\end{lemma}

\begin{proof}
By \eqref{eq:Sdef}, $w\in\OSP(n,n-c)$ with $0\le c\le b$, so
$n-b\le\blocks(w)=n-c\le n$, and $\inv(w)=\ell-b-c\in[\ell-2b,\ell-b]$.
\end{proof}

\begin{lemma}[Classification of $\phi$ on $S_{n,\ell,b}$]\label{lem:classify}
Let $w\in S_{n,\ell,b}$, say $w\in\OSP(n,n-c)$ with $\inv(w)=\ell-b-c$,
$0\le c\le b$. Then:
\begin{enumerate}[\rm(i)]
\item if $w$ is type I and $c=b$ (equivalently $\inv(w)=\ell-2b$), then
$\phi(w)\notin S_{n,\ell,b}$;
\item if $w$ is type I and $c<b$, then $\phi(w)\in S_{n,\ell,b}$, lying in
$\OSP(n,n-(c+1))$ with $\inv(\phi(w))=\ell-b-(c+1)$;
\item if $w$ is type II, then $c>0$ and $\phi(w)\in S_{n,\ell,b}$, lying in
$\OSP(n,n-(c-1))$ with $\inv(\phi(w))=\ell-b-(c-1)$;
\item $w$ cannot be type III.
\end{enumerate}
In cases (ii) and (iii), $\varepsilon(\phi(w))=-\varepsilon(w)$.
\end{lemma}

\begin{proof}
\emph{(iv)} Suppose $w$ were type III, i.e.\ a fixed point of $\phi$. A type III
ordered set partition has every block a singleton arranged in increasing order of
maxima, hence has $n$ blocks and $\inv=0$; in particular $\blocks(w)=n$, forcing
$c=0$, and $\inv(w)=0$, forcing $\ell-b=0$, i.e.\ $\ell=b$. But $\ell\ge1$ gives
$b=\ell>\lfloor(\ell-1)/2\rfloor$, contradicting $b\le\lfloor(\ell-1)/2\rfloor$.
So $w$ is type I or II.

\emph{(i)--(ii) (type I).} By Lemma~\ref{lem:phi-props}, $\phi(w)$ is type II with
$\blocks(\phi(w))=\blocks(w)-1=n-(c+1)$ and $\inv(\phi(w))=\inv(w)-1=\ell-b-c-1$.
For $\phi(w)$ to lie in $S_{n,\ell,b}$ we need its block count $n-(c+1)$ to be a
valid stratum, i.e.\ $0\le c+1\le b$, and its inversion number to match that
stratum, i.e.\ $\inv(\phi(w))=\ell-b-(c+1)$. The inversion condition holds
identically: $\ell-b-c-1=\ell-b-(c+1)$. Thus membership is governed solely by
$c+1\le b$:
\begin{itemize}
\item if $c=b$, then $c+1=b+1>b$, so $\phi(w)\notin S_{n,\ell,b}$; this is
(i), and indeed $\inv(w)=\ell-b-b=\ell-2b$.
\item if $c<b$, then $0\le c+1\le b$, so $\phi(w)\in\OSP(n,n-(c+1))$ with
$\inv(\phi(w))=\ell-b-(c+1)$, hence $\phi(w)\in S_{n,\ell,b}$; this is (ii).
\end{itemize}

\emph{(iii) (type II).} By Lemma~\ref{lem:phi-props}, $\phi(w)$ is type I with
$\blocks(\phi(w))=\blocks(w)+1=n-(c-1)$ and $\inv(\phi(w))=\inv(w)+1=\ell-b-c+1$.
Since $w$ is splittable, it has a block of size $\ge2$, so $\blocks(w)=n-c<n$,
forcing $c\ge1$, i.e.\ $c>0$; in particular $0\le c-1\le b$. The inversion
condition for the stratum $\OSP(n,n-(c-1))$ requires
$\inv(\phi(w))=\ell-b-(c-1)$, which holds: $\ell-b-c+1=\ell-b-(c-1)$. Hence
$\phi(w)\in S_{n,\ell,b}$.

\emph{Sign.} In cases (ii) and (iii), $\blocks(\phi(w))=\blocks(w)\pm1$, so by
\eqref{eq:epsdef} $\varepsilon(\phi(w))=(-1)^{\pm1}\varepsilon(w)=-\varepsilon(w)$.
\end{proof}

\begin{lemma}[Survivors]\label{lem:survivors}
Let
\begin{equation}\label{eq:Tdef}
T_{n,\ell,b}=\{w\in\OSP(n,n-b)\ :\ w\text{ is type I and }\inv(w)=\ell-2b\}.
\end{equation}
Then $\phi$ restricts to a fixed-point-free, sign-reversing involution on
$S_{n,\ell,b}\setminus T_{n,\ell,b}$, and $\phi$ maps $T_{n,\ell,b}$ out of
$S_{n,\ell,b}$. Consequently
\begin{equation}\label{eq:signed-sum}
\sum_{w\in S_{n,\ell,b}}\varepsilon(w) = \#\,T_{n,\ell,b}.
\end{equation}
\end{lemma}

\begin{proof}
First, $T_{n,\ell,b}\subseteq S_{n,\ell,b}$: an element of $T_{n,\ell,b}$ lies in
$\OSP(n,n-b)$ (so $c=b$) with $\inv=\ell-2b=\ell-b-b$, matching \eqref{eq:Sdef}.
By Lemma~\ref{lem:classify}(i) every $w\in T_{n,\ell,b}$ (type I, $c=b$) has
$\phi(w)\notin S_{n,\ell,b}$.

Now take $w\in S_{n,\ell,b}\setminus T_{n,\ell,b}$. By
Lemma~\ref{lem:classify}(iv) $w$ is type I or II. If $w$ is type I, then $w\notin
T_{n,\ell,b}$ forces $c<b$ (since type-I elements with $c=b$ are exactly
$T_{n,\ell,b}$), so case (ii) applies; if $w$ is type II, case (iii) applies. In
both cases $\phi(w)\in S_{n,\ell,b}$ and $\varepsilon(\phi(w))=-\varepsilon(w)$.

We check $\phi(w)\in S_{n,\ell,b}\setminus T_{n,\ell,b}$ as well, so that $\phi$
genuinely restricts to this subset. In case (ii) ($w$ type I, $c<b$),
Lemma~\ref{lem:phi-props} gives $\phi(w)$ of type II, and type II elements are
never in $T_{n,\ell,b}$ (which consists of type I elements); so
$\phi(w)\in S_{n,\ell,b}\setminus T_{n,\ell,b}$. In case (iii) ($w$ type II,
$c>0$), $\phi(w)$ is type I in $\OSP(n,n-(c-1))$; it lies in $T_{n,\ell,b}$ only
if its block count is $n-b$, i.e.\ $c-1=b$, i.e.\ $c=b+1$, impossible since
$c\le b$; so again $\phi(w)\in S_{n,\ell,b}\setminus T_{n,\ell,b}$.

Since $\phi$ is an involution on all ordered set partitions
(Lemma~\ref{lem:phi-props}) and maps
$S_{n,\ell,b}\setminus T_{n,\ell,b}$ into itself, it restricts to an involution
there. It is fixed-point-free on this subset: a fixed point of $\phi$ is type III,
excluded by Lemma~\ref{lem:classify}(iv). As it reverses $\varepsilon$, the
contributions of $S_{n,\ell,b}\setminus T_{n,\ell,b}$ to the signed sum cancel in
pairs:
\begin{equation}
\sum_{w\in S_{n,\ell,b}}\varepsilon(w)
=\sum_{w\in T_{n,\ell,b}}\varepsilon(w)
+\sum_{w\in S_{n,\ell,b}\setminus T_{n,\ell,b}}\varepsilon(w)
=\sum_{w\in T_{n,\ell,b}}\varepsilon(w).
\end{equation}
Finally every $w\in T_{n,\ell,b}$ has $\blocks(w)=n-b$, so $c=b$ and
$\varepsilon(w)=(-1)^{(n-b)-(n-b)}=(-1)^0=+1$ by \eqref{eq:epsdef}. Hence
$\sum_{w\in S_{n,\ell,b}}\varepsilon(w)=\#\,T_{n,\ell,b}$.
\end{proof}

\subsection{The main theorem}

\begin{theorem}\label{thm:improved-coefficient-formula}
For $a\ge1$ and $b\ge0$, the hook coefficient is given, manifestly positively, by
\begin{equation}\label{eq:second-formula-coeffs}
c_{(a,1^b)}(n)=\#\{w\in\OSP(n,n-b)\ :\ \inv(w)=a,\ w\text{ is type I}\}.
\end{equation}
Moreover $c_\varnothing(n)=1$ and $c_\lambda(n)=0$ for every non-hook $\lambda$.
\end{theorem}

\begin{proof}
Set $\ell:=a+2b$. Then $\ell\ge1$ (as $a\ge1$) and
$b\le\lfloor(\ell-1)/2\rfloor$, because $2b=\ell-a\le\ell-1$ gives
$b\le(\ell-1)/2$ and $b$ is an integer. With this $\ell$, the defining set
\eqref{eq:Sdef} becomes
\begin{equation}
S_{n,a+2b,b}=\bigcup_{c=0}^b\{w\in\OSP(n,n-c):\inv(w)=a+b-c\}.
\end{equation}
Reindex the alternating formula \eqref{eq:first-formula-coeffs} by
$c:=b-i$ (so $i=b-c$ runs over $0,\dots,b$ as $c$ runs over $b,\dots,0$). Then
$\OSP(n,n-b+i)=\OSP(n,n-c)$, $\inv=a+i=a+b-c$, and the sign is
$(-1)^i=(-1)^{b-c}=\varepsilon(w)$ for $w\in\OSP(n,n-c)$, by \eqref{eq:epsdef}.
Therefore
\begin{equation}
c_{(a,1^b)}(n)
=\sum_{c=0}^b(-1)^{b-c}\#\{w\in\OSP(n,n-c):\inv(w)=a+b-c\}
=\sum_{w\in S_{n,a+2b,b}}\varepsilon(w).
\end{equation}
By Lemma~\ref{lem:survivors} with $\ell=a+2b$,
\begin{equation}
c_{(a,1^b)}(n)=\#\,T_{n,a+2b,b}
=\#\{w\in\OSP(n,n-b):w\text{ type I},\ \inv(w)=(a+2b)-2b\},
\end{equation}
and $(a+2b)-2b=a$, which is exactly \eqref{eq:second-formula-coeffs}. This count
is a cardinality, hence a nonnegative integer, so the formula is manifestly
positive.

The statement $c_\varnothing(n)=1$ is the coefficient of $q^0u^0$ in
\eqref{eq:master}: the constant term on the left is $1$ and on the right is the
$d=n$, $i=0$ contribution $\#\{w\in\OSP(n,n):\inv(w)=0\}=1$ (the unique increasing
ordered set partition into singletons), confirming $c_\varnothing(n)=1$. For a
non-hook $\lambda$ ($\lambda_2\ge2$) we have $\lambda\notin P(1,1,n)$, so it does
not occur in \eqref{eq:coeffs-Hilbert}; equivalently $s_\lambda(q/u)=0$ by
Lemma~\ref{lem:schur-1-1}(iii). By the universality in
Proposition~\ref{prop:bergeron} and the supersymmetric identification in
Proposition~\ref{prop:ds}, the coefficient $c_\lambda(n)$ of every non-hook
$\lambda$ is $0$.
\end{proof}

\section{Example: $n=4$}
The following table lists the nonzero hook coefficients $c_\lambda(4)$ together
with the ordered set partitions counted by
Theorem~\ref{thm:improved-coefficient-formula}. (All blocks are written
increasingly; e.g.\ $2|1|3|4$ is $(\{2\}/\{1\}/\{3\}/\{4\})$.)

\begin{figure}[h]
\begin{tabular}{l|l|llllll}
\text{Coefficient} & \text{Value} & \text{OSPs}         \\ \hline
$c_{(1)}(4)$       & 3            & $2|1|3|4$  & $1|3|2|4$ & $1|2|4|3$ &           &           &           \\
$c_{(1,1)}(4)$     & 3            & $1|4|23$   & $12|4|3$  & $3|12|4$  &           &           &           \\
$c_{(1,1,1)}(4)$   & 1            & $4|123$    &           &           &           &           &           \\
$c_{(2)}(4)$       & 5            & $2|1|4|3$  & $1|3|4|2$ & $2|3|1|4$ & $3|1|2|4$ & $1|4|2|3$ &           \\
$c_{(2,1)}(4)$     & 5            & $4|1|23$   & $13|4|2$  & $2|4|13$  & $3|4|12$  & $4|12|3$  &           \\
$c_{(3)}(4)$       & 6            & $2|3|4|1$  & $3|2|1|4$ & $2|4|1|3$ & $3|1|4|2$ & $4|1|2|3$ & $1|4|3|2$ \\
$c_{(3,1)}(4)$     & 4            & $4|13|2$   & $23|4|1$  & $4|2|13$  & $4|3|12$  &           &           \\
$c_{(4)}(4)$       & 5            & $4|1|3|2$  & $4|2|1|3$ & $3|4|1|2$ & $3|2|4|1$ & $2|4|3|1$ &           \\
$c_{(4,1)}(4)$     & 1            & $4|23|1$   &           &           &           &           &           \\
$c_{(5)}(4)$       & 3            & $4|2|3|1$  & $4|3|1|2$ & $3|4|2|1$ &           &           &           \\
$c_{(6)}(4)$       & 1            & $4|3|2|1$  &           &           &           &           &          
\end{tabular}
\caption{The nonzero hook coefficients $c_\lambda(4)$ and the ordered set
partitions counted by Theorem~\ref{thm:improved-coefficient-formula}. Apart from
$c_\varnothing(4)=1$, all other $c_\lambda(4)$ are $0$.}
\end{figure}

These reproduce the Schur (super Schur) expansion
\begin{align*}
\Hilb(R_4^{(1,1)};q;u)
&= (q^6+3q^5+5q^4+6q^3+5q^2+3q+1)
 + (q^5+4q^4+9q^3+11q^2+8q+3)u\\
&\quad + (q^3+4q^2+6q+3)u^2 + u^3\\
&= 1+\sum_{a=1}^{6} c_{(a)}(4)\,s_{(a)}(q/u)
 + \sum_{a=1}^{4} c_{(a,1)}(4)\,s_{(a,1)}(q/u)
 + c_{(1,1,1)}(4)\,s_{(1,1,1)}(q/u),
\end{align*}
using $s_{(a,1^b)}(q/u)=q^au^b+q^{a-1}u^{b+1}$ from
Lemma~\ref{lem:schur-1-1}.

\begin{thebibliography}{99}
\bibitem{Bergeron2020} Fran\c{c}ois Bergeron, \emph{The bosonic-fermionic
diagonal coinvariant modules conjecture}, preprint (2020),
\url{https://arxiv.org/abs/2005.00924}.
\bibitem{ChengWangBook} Shun-Jen Cheng and Weiqiang Wang, \emph{Dualities and
representations of Lie superalgebras}, Graduate Studies in Mathematics, vol.\
144, American Mathematical Society, Providence, RI, 2012.
\bibitem{Lentfer-Supersymmetry} John Lentfer, \emph{Diagonal supersymmetry for
coinvariant rings}, preprint (2025), \url{https://arxiv.org/abs/2505.14885}.
\bibitem{RhoadesWilson2023} Brendon Rhoades and Andrew Timothy Wilson, \emph{The
Hilbert series of the superspace coinvariant ring}, Forum Math.\ Pi 12 (2024),
Id/No e16.
\bibitem{SaganSwanson2024} Bruce E. Sagan and Joshua P. Swanson, \emph{$q$-Stirling
numbers in type $B$}, Eur.\ J.\ Comb.\ 118 (2024), Id/No 103899.
\end{thebibliography}

\end{document}
